\section{The constants by interval arithmetic}
\label{app:constants}

No step of the proof is computer-assisted: every constant is a closed form.
To exclude arithmetic slips we recomputed every one of them with a small interval arithmetic on Python's \code{decimal} module (\code{src/eh\_iv.py}): $60$ digits, the lower end rounded down and the upper end rounded up in every $+,-,\times,/$, and $\sqrt\cdot$, $\exp$, $\ln$ correctly rounded and then widened by one unit in the last place on each side; $\pi$ enters as an interval from its decimal expansion.
A claim holds when the whole enclosure of the exact expression lies on the stated side of the value used in the text.
Claims for all $L\ge4$ or all $\sigma\le\sigma_0$ are proved in the text by monotonicity; the table checks their end points, and \code{src/eh\_constants.py} also samples them on logarithmic grids ($L=4$ to $2\times10^7$, $\sigma=10^{-15}$ to $10^{-300}$).
All $101$ claims of Table~\ref{tab:constants} hold.

The same check of the rounded values printed in the working record found three that were rounded in the wrong direction for the way they were used (Table~\ref{tab:printed}).
Two of them enter this paper, in Theorem~\ref{thm:circles} ($\sqrt{4/\ln8}=1.38693\ldots$, printed as $1.387$ and used as a lower bound) and Lemma~\ref{lem:iso} ($4/\ln8=1.92359\ldots$, printed as $1.924$); the text now uses the correct values, and the constants $36.1$ and $70.2$ derived from them hold unchanged.
The third belongs to a lemma that the chain does not use.

% generated by src/ei_tables.py from results_EH/eh_constants.json (the strings of the record set as formulas by PRINTED) -- do not edit by hand
\begin{table}[t]
\caption{Rounded values printed in the working record and used there as one-sided bounds, checked for their direction against the enclosure of the exact value. The three marked ``wrong'' are corrected in this paper; the constants derived from them hold when recomputed (Table~\ref{tab:constants}). The last two were used as lower bounds. $^\dagger$The original form of Lemma EE-C, not used by Theorem~\ref{thm:main}.}
\label{tab:printed}
\begin{ruledtabular}
\footnotesize\setlength{\tabcolsep}{2pt}
\begin{tabular}{@{}llll@{}}
record & printed value & enclosure & check\\
\hline
EE 3.3$^\dagger$ & $c_1=2\cdot315^{1/4}+\sqrt{15}/3\le9.7167$ & 9.71673 & \textbf{wrong}\\
EG 2.4 step 7 & $1-\tau\ge0.9837$ at $L=4$ & 0.983754 & right\\
EG 4.2 step 4 & $(1-\tau)^3\ge0.952$ at $L=4$ & 0.952051 & right\\
EF 2.4 & $1.291\sqrt{112}\le13.67$ & 13.6627 & right\\
EF 2.4 & $8.4258\sqrt{224}\le126.11$ & 126.106 & right\\
EF 2.4 & $4\sqrt{224\cdot112}\le633.6$ & 633.568 & right\\
EG 4.3 & $\alpha\ge0.0506$ & 0.0506606 & right\\
EG 4.3 & $\Xi^2/\alpha\le0.0132$ & 0.0131225 & right\\
EG 4.3 & $\Xi/(2\sqrt\alpha)\le0.0574$ & 0.0573475 & right\\
EG 2.4 step 7 & $\sqrt{4/\ln8}\ge1.387$ & 1.38694 & \textbf{wrong}\\
EG 4.2 step 4 & $4/\ln8\ge1.924$ & 1.92359 & \textbf{wrong}\\
\end{tabular}
\end{ruledtabular}
\end{table}


% generated by src/ei_tables.py from results_EH/eh_constants.json and results_EI/ei_constants.json -- do not edit by hand
% Table 4 in 3 groups of at most 34 rows (plain table* floats)
\begin{table*}[p]
\scriptsize
\caption{Every constant of the chain, recomputed with outward-rounded interval arithmetic (60 digits; \code{src/eh\_constants.py}, \code{src/ei\_constants.py}). The expressions are printed as evaluated by the code; two explanatory remarks of the record are reworded. The enclosure is the interval of the exact expression; ``claimed'' is the value used in the text, and ``side'' says on which side the enclosure must lie. All 101 hold. The table continues in 2 more parts.}\label{tab:constants}
\begin{ruledtabular}
\begin{tabular}{@{}lp{0.50\textwidth}lll@{}}
part & claim (the expression, as evaluated by the code) & lower end of the enclosure & claimed & side\\
\hline
L.\ref{lem:CR} & \texttt{alpha = 1/(2(pi\textasciicircum{}2 + sigma\textasciicircum{}2)) at sigma = 1e-4 (alpha decreases in sigma)} & 0.0506606 & 0.0506 & $\ge$\\
L.\ref{lem:K} & \texttt{C\_\{2,1/4\} = M\_2\textasciicircum{}\{1/2\} M\_2\textasciicircum{}\{1/2\} (1/2)\textasciicircum{}\{-5/4\} = 3 * 2\textasciicircum{}\{5/4\}} & 7.13524 & 7.136 & $\le$\\
L.\ref{lem:K} & \texttt{coefficient r'/lambda = 2/(1/4)} & 8 & 8 & $\le$\\
L.\ref{lem:G2} & \texttt{(sqrt960)/6  (E delta\textasciicircum{}6 $\le$ 2\textasciicircum{}6 * 15 sigma\textasciicircum{}6)} & 5.16398 & 5.164 & $\le$\\
L.\ref{lem:EF1} & \texttt{\textbar{}\textbar{}X\_+-\textbar{}\textbar{}\_2 $\le$ 2 sqrt2 + sqrt(3/2)} & 4.05317 & 4.0533 & $\le$\\
L.\ref{lem:EF1} & \texttt{c1 = sqrt3/2 + sqrt3 + sqrt960/6} & 7.76205 & 7.77 & $\le$\\
L.\ref{lem:EF1} & \texttt{c2 = (2 sqrt2 + sqrt(3/2))/sqrt2} & 2.86603 & 2.87 & $\le$\\
L.\ref{lem:EF1} & \texttt{c3 = (2 sqrt2 + sqrt(3/2))/2} & 2.02659 & 2.03 & $\le$\\
L.\ref{lem:EF2} & \texttt{c4 = sqrt(2.03\textasciicircum{}2 + 2 * 2.87\textasciicircum{}2)} & 4.53814 & 4.5382 & $\le$\\
L.\ref{lem:EF2} & \texttt{(1/2) 112 - c4 sqrt112 - 7.77  ($>$ 0: the threshold M $\ge$ 112 sigma\textasciicircum{}2)} & 0.202858 & 0 & $>$\\
L.\ref{lem:EF2} & \texttt{alternative: 4 (c1 + c2\textasciicircum{}2 + c3\textasciicircum{}2)  (A $\ge$ 80.5 sigma\textasciicircum{}2 suffices)} & 80.3328 & 80.5 & $\le$\\
L.\ref{lem:EEC} & \texttt{sqrt15/3  (\textbar{}E[s\_v\textasciicircum{}3 s\_w]\textbar{}/3 $\le$ sqrt(15 A)/3)} & 1.29099 & 1.291 & $\le$\\
L.\ref{lem:EEC} & \texttt{c\_2 = sqrt(2 sqrt1260)  (E G\textasciicircum{}4 $\le$ sqrt(105 * 1680) = 420)} & 8.42573 & 8.4258 & $\le$\\
L.\ref{lem:EEC} & \texttt{105 * 1680 = 420\textasciicircum{}2 (integers)} & 176400 & 176400 & $=$\\
L.\ref{lem:EFO} & \texttt{1.291 sqrt112 + 8.4258 sqrt224 + 112 sigma  (sigma $\le$ 1e-3)} & 139.88 & 140 & $\le$\\
L.\ref{lem:EFO} & \texttt{8 sqrt(224 * 112)  (B $\le$ 2 * 4 k sqrt(theta B A)/sigma\textasciicircum{}2 ... ; C\textasciitilde{}\_w $\ge$ 1/2)} & 1267.14 & 1268 & $\le$\\
P.\ref{prop:outside} & \texttt{beta' at sigma = 1e-4 (decreasing as sigma decreases)} & 0.283818 & 0.284 & $\le$\\
P.\ref{prop:outside} & \texttt{112/K = 112/400} & 0.28 & 0.28 & $\le$\\
T.\ref{thm:circles} & \texttt{eta = ln(2L)/(4L) at L = 4 (decreasing)} & 0.129965 & 0.13 & $<$\\
T.\ref{thm:circles} & \texttt{s = ln(2L)/(8L\textasciicircum{}2) at L = 4 (decreasing)} & 0.0162456 & 0.0163 & $<$\\
T.\ref{thm:circles} & \texttt{-ln(1 - eta) - 2 eta at eta = 0.13 ($\le$ 0; it is $\le$ 0 for every eta in [0, 1/2], since it vanishes at 0 and its derivative 1/(1 - eta) - 2 is negative there)} & -0.120738 & 0 & $\le$\\
T.\ref{thm:circles} & \texttt{r2 $<$ r3 at L = 4} & 0.0787456 & 0.53125 & $<$\\
T.\ref{thm:circles} & \texttt{theta = ln(r2/r1)/ln(r3/r1) at L = 4, times L ($<$ 1)} & 0.649814 & 1 & $<$\\
T.\ref{thm:circles} & \texttt{alpha K\textasciicircum{}2 = 4 L alpha at L = 4 ($<$ 1)} & 0.640368 & 1 & $<$\\
T.\ref{thm:circles} & \texttt{T/(1-s) at L = 4} & 1.08005 & 1.0801 & $\le$\\
T.\ref{thm:circles} & \texttt{sqrt(6 T/(1 - s))  (contribution of E\_bad, x A)} & 2.54564 & 2.55 & $\le$\\
T.\ref{thm:circles} & \texttt{4e/sqrt3  (three circles, x sqrt(L/ln 2L) A)} & 6.2776 & 6.2777 & $\le$\\
T.\ref{thm:circles} & \texttt{2e * 4e/sqrt3  (outside E\_bad)} & 34.1286 & 34.13 & $\le$\\
T.\ref{thm:circles} & \texttt{sqrt(4/ln 8) $\ge$ 1.3869 (the correct lower bound)} & 1.38694 & 1.3869 & $\ge$\\
T.\ref{thm:circles} & \texttt{total: 2e (4e/sqrt3) + 2.55/sqrt(4/ln8)} & 35.9672 & 36.1 & $\le$\\
T.\ref{thm:circles} & \texttt{36.1 * 8  (Q2 $\le$ F(1-s)/s, 1/s = 8L\textasciicircum{}2/ln 2L)} & 288.8 & 289 & $\le$\\
T.\ref{thm:circles} & \texttt{psi(4) = 289 * 32/(ln 8)\textasciicircum{}\{3/2\} + 19.57} & 3103.67 & 3100 & $\ge$\\
T.\ref{thm:circles} & \texttt{e\textasciicircum{}4 + 4*4 + 3.57  (A $>$ e\textasciicircum{}-4: k\_eff\textasciicircum{}2 $<$ e\textasciicircum{}4 + 4L + 3.57)} & 74.1682 & 74.2 & $\le$\\
L.\ref{lem:iso},\ref{lem:sm} & \texttt{outer, isotropic: T\textasciicircum{}3/(T\textasciicircum{}2 - (T/2)\textasciicircum{}2)\textasciicircum{}2 = 16/(9T) $\le$ 16/9} & 1.6732 & 1.77778 & $\le$\\
\end{tabular}
\end{ruledtabular}
\end{table*}

\begin{table*}[p]
\scriptsize
{\small\centering Table~\ref{tab:constants} (continued)\par}
\smallskip
\begin{ruledtabular}
\begin{tabular}{@{}lp{0.50\textwidth}lll@{}}
part & claim (the expression, as evaluated by the code) & lower end of the enclosure & claimed & side\\
\hline
L.\ref{lem:iso},\ref{lem:sm} & \texttt{three circles, isotropic: e (16/9)\textasciicircum{}\{1/4\} * 4} & 12.5552 & 12.56 & $\le$\\
L.\ref{lem:iso},\ref{lem:sm} & \texttt{sqrt6 T\textasciicircum{}2/(1 - s)\textasciicircum{}3  (E\_bad, isotropic)} & 2.90451 & 2.91 & $\le$\\
L.\ref{lem:iso},\ref{lem:sm} & \texttt{4/ln 8 $\ge$ 1.9235 (the correct lower bound)} & 1.92359 & 1.9235 & $\ge$\\
L.\ref{lem:iso},\ref{lem:sm} & \texttt{total: 2e * 12.56 + 2.91/(4/ln 8)} & 69.796 & 70.2 & $\le$\\
L.\ref{lem:iso},\ref{lem:sm} & \texttt{70.2 * 8  (E[(d\_e\textasciicircum{}2 p/p)\textasciicircum{}2] $\le$ G(1-s)/s)} & 561.6 & 562 & $\le$\\
L.\ref{lem:EGD} & \texttt{sqrt15/6  (\textbar{}R'\textbar{} $\le$ (sqrt15/6) sigma\textasciicircum{}2 sqrt(J\_pp))} & 0.645497 & 0.6455 & $\le$\\
T.\ref{thm:EGO} & \texttt{h(2e-9) = 562 J0 L\textasciicircum{}3/(ln 2L)\textasciicircum{}2} & 0.000663245 & 0.000664 & $\le$\\
T.\ref{thm:EGO} & \texttt{Xi = sqrt(h(2e-9))} & 0.0257535 & 0.0258 & $\le$\\
T.\ref{thm:EGO} & \texttt{bracket Xi\textasciicircum{}2/alpha + Xi/(2 sqrt alpha) + 0.6455 sqrt(J0)/alpha} & 0.0708717 & 0.25 & $\le$\\
T.\ref{thm:EGO} & \texttt{h increasing on (0, e\textasciicircum{}-4]: 1 - 3/L + 2/(L ln 2L) at L = 4} & 0.490449 & 0 & $>$\\
T.\ref{thm:main} & \texttt{L\_bar = 2 ln(1e15) + ln(2(pi\textasciicircum{}2 + sigma\textasciicircum{}2))} & 72.0602 & 72 & $\ge$\\
T.\ref{thm:main} & \texttt{k\_bar = sqrt(psi(L\_bar)) at sigma = 1e-15} & 1072.29 & 1072.4 & $\le$\\
T.\ref{thm:main} & \texttt{J inside O $\le$ 112 sigma\textasciicircum{}2 + 1268 sigma k\_bar at sigma = 1e-15} & 1.35967e-9 & 2e-9 & $\le$\\
T.\ref{thm:main} & \texttt{(3/4) alpha/(alpha sigma + 1268 k\_bar) * k\_bar} & 2.99649e-5 & 2.9e-5 & $\ge$\\
T.\ref{thm:main} & \texttt{psi(L) $\le$ 289.1 L\textasciicircum{}\{5/2\}/(ln 2L)\textasciicircum{}\{3/2\} at L = 72: (4L + 3.57)(ln 2L)\textasciicircum{}\{3/2\}/L\textasciicircum{}\{5/2\}} & 0.0734381 & 0.1 & $\le$\\
T.\ref{thm:main} & \texttt{sqrt(289.1)} & 17.0029 & 17.01 & $\le$\\
T.\ref{thm:main} & \texttt{0.75 alpha/1268/17.01  ($\ge$ 1.7e-6)} & 1.7616e-6 & 1.7e-6 & $\ge$\\
T.\ref{thm:main} & \texttt{bound $\ge$ 1.7e-6 (ln 2L)\textasciicircum{}\{3/4\}/L\textasciicircum{}\{5/4\} at sigma = 1e-15} & 2.79447e-8 & 2.6955e-8 & $\ge$\\
T.\ref{thm:main} & \texttt{J inside O $\le$ 2e-9 at sigma = 1e-15} & 1.35967e-9 & 2e-9 & $\le$\\
T.\ref{thm:main} & \texttt{the O-term is the smallest of the five at sigma = 1e-15} & 2.79447e-8 & 0.000118039 & $<$\\
T.\ref{thm:main} & \texttt{bound $\ge$ 1.7e-6 (ln 2L)\textasciicircum{}\{3/4\}/L\textasciicircum{}\{5/4\} at sigma = 1e-20} & 2.05811e-8 & 1.98515e-8 & $\ge$\\
T.\ref{thm:main} & \texttt{J inside O $\le$ 2e-9 at sigma = 1e-20} & 1.84613e-14 & 2e-9 & $\le$\\
T.\ref{thm:main} & \texttt{the O-term is the smallest of the five at sigma = 1e-20} & 2.05811e-8 & 8.69351e-5 & $<$\\
T.\ref{thm:main} & \texttt{bound $\ge$ 1.7e-6 (ln 2L)\textasciicircum{}\{3/4\}/L\textasciicircum{}\{5/4\} at sigma = 1e-30} & 1.32651e-8 & 1.27944e-8 & $\ge$\\
T.\ref{thm:main} & \texttt{J inside O $\le$ 2e-9 at sigma = 1e-30} & 2.86431e-24 & 2e-9 & $\le$\\
T.\ref{thm:main} & \texttt{the O-term is the smallest of the five at sigma = 1e-30} & 1.32651e-8 & 5.60321e-5 & $<$\\
T.\ref{thm:main} & \texttt{bound $\ge$ 1.7e-6 (ln 2L)\textasciicircum{}\{3/4\}/L\textasciicircum{}\{5/4\} at sigma = 1e-50} & 7.5474e-9 & 7.2794e-9 & $\ge$\\
T.\ref{thm:main} & \texttt{J inside O $\le$ 2e-9 at sigma = 1e-50} & 5.03425e-44 & 2e-9 & $\le$\\
T.\ref{thm:main} & \texttt{the O-term is the smallest of the five at sigma = 1e-50} & 7.5474e-9 & 3.18803e-5 & $<$\\
T.\ref{thm:main} & \texttt{bound $\ge$ 1.7e-6 (ln 2L)\textasciicircum{}\{3/4\}/L\textasciicircum{}\{5/4\} at sigma = 1e-100} & 3.46342e-9 & 3.34038e-9 & $\ge$\\
T.\ref{thm:main} & \texttt{J inside O $\le$ 2e-9 at sigma = 1e-100} & 1.09705e-93 & 2e-9 & $\le$\\
T.\ref{thm:main} & \texttt{the O-term is the smallest of the five at sigma = 1e-100} & 3.46342e-9 & 1.46295e-5 & $<$\\
T.\ref{thm:main} & \texttt{bound $\ge$ 1.7e-6 (ln 2L)\textasciicircum{}\{3/4\}/L\textasciicircum{}\{5/4\} at sigma = 1e-300} & 9.85903e-10 & 9.5087e-10 & $\ge$\\
T.\ref{thm:main} & \texttt{J inside O $\le$ 2e-9 at sigma = 1e-300} & 3.85387e-293 & 2e-9 & $\le$\\
\end{tabular}
\end{ruledtabular}
\end{table*}

\begin{table*}[p]
\scriptsize
{\small\centering Table~\ref{tab:constants} (continued)\par}
\smallskip
\begin{ruledtabular}
\begin{tabular}{@{}lp{0.50\textwidth}lll@{}}
part & claim (the expression, as evaluated by the code) & lower end of the enclosure & claimed & side\\
\hline
T.\ref{thm:main} & \texttt{the O-term is the smallest of the five at sigma = 1e-300} & 9.85903e-10 & 4.16447e-6 & $<$\\
T.\ref{thm:EGO},\ref{thm:main} & \texttt{bracket of Theorem EG-O at J0 = 1.9e-8} & 0.247013 & 0.25 & $\le$\\
T.\ref{thm:EGO},\ref{thm:main} & \texttt{J inside O $\le$ 112 sigma\textasciicircum{}2 + 1268 sigma k\_bar at sigma = 1.5e-14} & 1.87186e-8 & 1.9e-8 & $\le$\\
T.\ref{thm:EGO},\ref{thm:main} & \texttt{alpha $\ge$ 0.050660 for sigma\textasciicircum{}2 $\le$ J0/alpha (sigma\textasciicircum{}2 $\le$ 4e-7)} & 0.0506606 & 0.05066 & $\ge$\\
T.\ref{thm:EGO} & \texttt{sigma\textasciicircum{}2 $\le$ 2 pi\textasciicircum{}2 J0/(1 - 2 J0)} & 3.75045e-7 & 3.751e-7 & $\le$\\
T.\ref{thm:EGO} & \texttt{alpha $\ge$ 1/(2(pi\textasciicircum{}2 + sigma\textasciicircum{}2\_max))} & 0.0506606 & 0.0506605 & $\ge$\\
T.\ref{thm:EGO} & \texttt{L\_J = ln(1/J0) $\ge$ 4 (h\_O increasing there)} & 17.7788 & 4 & $\ge$\\
T.\ref{thm:EGO} & \texttt{h\_O(J0) = 562 J0 L\textasciicircum{}3/(ln 2L)\textasciicircum{}2} & 0.00470524 & 0.005105 & $\le$\\
T.\ref{thm:EGO} & \texttt{Xi = sqrt(h\_O(J0))} & 0.0685947 & 0.07145 & $\le$\\
T.\ref{thm:EGO} & \texttt{Xi\textasciicircum{}2/alpha} & 0.0928777 & 0.1008 & $\le$\\
T.\ref{thm:EGO} & \texttt{Xi/(2 sqrt alpha)} & 0.152379 & 0.1588 & $\le$\\
T.\ref{thm:EGO} & \texttt{(sqrt15/6) sqrt(J0)/alpha} & 0.00175631 & 0.00176 & $\le$\\
T.\ref{thm:EGO} & \texttt{the bracket Xi\textasciicircum{}2/alpha + Xi/(2 sqrt alpha) + (sqrt15/6) sqrt(J0)/alpha} & 0.247013 & 0.2471 & $\le$\\
T.\ref{thm:EGO} & \texttt{the bracket $\le$ 1/4} & 0.247013 & 0.25 & $\le$\\
T.\ref{thm:main} & \texttt{L\_bar at sigma = 1.5e-14} & 66.6441 & 66.6 & $\ge$\\
T.\ref{thm:main} & \texttt{k\_bar at sigma = 1.5e-14} & 984.152 & 984.16 & $\le$\\
T.\ref{thm:main} & \texttt{A + B inside O $\le$ 112 sigma\textasciicircum{}2 + 1268 sigma k\_bar at sigma = 1.5e-14} & 1.87186e-8 & 1.872e-8 & $\le$\\
T.\ref{thm:main} & \texttt{that is $\le$ J0 = 1.9e-8} & 1.87186e-8 & 1.9e-8 & $\le$\\
T.\ref{thm:main} & \texttt{beta' $\le$ 0.284 at sigma = 1.5e-14} & 0.232341 & 0.284 & $\le$\\
T.\ref{thm:main} & \texttt{the O-term (3/4) alpha/(alpha sigma + 1268 k\_bar) is the smallest of the five} & 3.04474e-8 & 0.00012861 & $<$\\
T.\ref{thm:main} & \texttt{the O-term at sigma = 1.5e-14} & 3.04474e-8 & 3.044e-8 & $\ge$\\
T.\ref{thm:main} & \texttt{the O-term times k\_bar} & 2.99649e-5 & 2.996e-5 & $\ge$\\
T.\ref{thm:main} & \texttt{L\_bar at sigma = 1e-15} & 72.0602 & 72.06 & $\ge$\\
T.\ref{thm:main} & \texttt{k\_bar at sigma = 1e-15} & 1072.29 & 1072.3 & $\le$\\
T.\ref{thm:main} & \texttt{A + B inside O $\le$ 112 sigma\textasciicircum{}2 + 1268 sigma k\_bar at sigma = 1e-15} & 1.35967e-9 & 1.3598e-9 & $\le$\\
T.\ref{thm:main} & \texttt{that is $\le$ J0 = 1.9e-8} & 1.35967e-9 & 1.9e-8 & $\le$\\
T.\ref{thm:main} & \texttt{beta' $\le$ 0.284 at sigma = 1e-15} & 0.229491 & 0.284 & $\le$\\
T.\ref{thm:main} & \texttt{the O-term (3/4) alpha/(alpha sigma + 1268 k\_bar) is the smallest of the five} & 2.79447e-8 & 0.000118039 & $<$\\
T.\ref{thm:main} & \texttt{the O-term at sigma = 1e-15} & 2.79447e-8 & 2.794e-8 & $\ge$\\
T.\ref{thm:main} & \texttt{the O-term times k\_bar} & 2.99649e-5 & 2.996e-5 & $\ge$\\
T.\ref{thm:main} & \texttt{1.7e-6 (ln 2 L\_bar)\textasciicircum{}\{3/4\}/L\_bar\textasciicircum{}\{5/4\} $\le$ the O-term at sigma = 1e-15} & 2.6955e-8 & 2.79447e-8 & $\le$\\
R.\ref{rem:delta} & \texttt{delta = 0.4/L at L = 4: omega K\_b\textasciicircum{}2 = 4 L omega = (0.8 + ln(2L)/(4L))/(1 - tau) $<$ 1} & 0.945322 & 0.9455 & $\le$\\
T.\ref{thm:main} & \texttt{psi'(L)/psi(L) * L/2.5 at L = 4 ($<$ 1)} & 0.709037 & 1 & $<$\\
\end{tabular}
\end{ruledtabular}
\end{table*}


